\documentclass[../../../include/open-logic-section]{subfiles}

\begin{document}

\olfileid[pt]{sth}{z}{arbintersections}
\olsection{Apêndice: Fechamento, compreensão e interseção}

Na \olref[sth][z][milestone]{sec}, sugerimos que voltasse atrás no trabalho
ingênuo da \olref[sfr][][]{part} e verificasse que este pode ser
realizado em~$\Zminus$. Se seguiu esse conselho,
\emph{um} ponto poderá tê-lo feito tropeçar: o uso da
\emph{interseção} no tratamento por Dedekind dos \emph{fechamentos}.

Recordemos, de \olref[sfr][infinite][dedekind]{Closure}, que
\begin{align*}
  \closureofunder{f}{o} & = \bigcap\Setabs{X}{o \in X \text{ e $X$ é
$f$-fechado}}.
\intertext{A forma geral disto é uma definição da forma:}
  C & = \bigcap\Setabs{X}{\phi(X)}.
\end{align*}
Mas isto deve fazer soar o alarme: uma vez que a compreensão ingênua
falha, não há garantia de que $\Setabs{X}{\phi(X)}$ exista. Parece,
então, perigosamente, que tais definições são \emph{truques}.

Felizmente, não são truques; ou melhor, se \emph{o forem}
tal como estão, então podemos empenhar-nos honestamente para
as tornar legítimas. Esse trabalho honesto foi antecipado na
\olref[sth][z][sep]{prop:intersectionsexist}, quando explicamos por que
$\bigcap A$ existe para qualquer $A \neq \emptyset$. Mas vamos explicitá-lo.

Dada a Extensionalidade, se tentarmos definir $C$ como
$\bigcap\Setabs{X}{\phi(X)}$, tudo o que realmente pedimos é um
objeto $C$ que satisfaça o seguinte:
\begin{equation}\ollabel{bicondelimarbintersection}
	\forall x(x \in C \liff \forall X(\phi(X) \lif x \in X))\tag{*}
\end{equation}
Suponhamos, agora, que existe \emph{algum} conjunto, $S$, tal que $\phi(S)$. Então,
para satisfazer \olref{bicondelimarbintersection}, podemos simplesmente definir $C$
por \emph{Separação}, do seguinte modo:
\[
	C = \Setabs{x \in S}{\forall X(\phi(X) \lif x \in X)}.
\]
Deixamos como exercício verificar que esta definição produz
\olref{bicondelimarbintersection}, como desejado. 
%  fix $x$. First suppose that $x \in C$, as (re)defined; then both $x
%  \in S$ and $\forall X(\phi(X) \lif x \in X)$; but since
%  $\phi(S)$, this reduces to the condition that $\forall X(\phi(X)
%  \lif x \in X)$. Conversely, suppose $\forall X(\phi(X) \lif
%  x \in X)$; then, since $\phi(S)$, we also have that $x \in S$; so
%  $x \in c$, as (re)defined. 
E esta estratégia geral permitir-nos-á contornar qualquer uso aparente da
compreensão ingênua na definição de interseções. No caso particular que nos
levou a esta linha de raciocínio, nomeadamente o de
$\closureofunder{f}{o}$, eis como isso funcionaria. Começamos a
prova de \olref[sfr][infinite][dedekind]{closureproperties} ao notar
que $o \in \ran{f}\cup\{o\}$ e que $\ran{f} \cup \{o\}$ é
$f$-fechado. Assim, podemos definir o que queremos assim:
\[
	\closureofunder{f}{o} = \Setabs{x \in \ran{f} \cup \{o\}}{(\forall X \ni o)(X \text{ é $f$-fechado} \lif x \in X)}.
\]

\end{document}
